\documentclass[11pt]{article}
\usepackage[T1]{fontenc}
\usepackage[utf8]{inputenc}
\usepackage{lmodern}
\usepackage{amsmath,amssymb}
\usepackage{longtable,booktabs,array}
% Compatibility with Pandoc's captionless tables on older longtable versions.
\ifcsname c@none\endcsname\else\newcounter{none}\fi
\usepackage{graphicx}
\usepackage[margin=25mm]{geometry}
\usepackage{hyperref}
\hypersetup{colorlinks=true,linkcolor=blue,urlcolor=blue}
\providecommand{\tightlist}{\setlength{\itemsep}{0pt}\setlength{\parskip}{0pt}}
\providecommand{\pandocbounded}[1]{#1}
\setlength{\emergencystretch}{3em}
\setcounter{secnumdepth}{-1}
\begin{document}
\section{Cohomology of projective
space}\label{cohomology-of-projective-space}

\emph{Written by GPT-6.1 Sol (OpenAI), in Codex, at Ultra effort,
October 2026. Self-checked by the writing AI, GPT-6.1 Sol, at Ultra
effort. Public domain (CC0).}

Projective space has a finite affine cover, so its cohomology can be
calculated from localizations of a polynomial ring. The calculation
separates according to individual Laurent monomials. Their negative
exponents determine which intersections can carry them, and this
elementary combinatorics explains both the vanishing of intermediate
cohomology and the duality between ordinary and reciprocal monomials.

We assume \href{affine-cohomology-and-serres-criterion.md}{Cohomology of
affine schemes and Serre's criterion}. We construct the projective
charts, twisting sheaves and quotient convention here, before using
them. No calculation requires a field, a reduced ring, or a
characteristic-zero hypothesis.

For clarity, the degree-one Proj construction is the following gluing.
If \(B\) is a graded \(A\)-algebra generated by degree-one elements
\(b_i\), glue \(\operatorname{Spec}(B[b_i^{-1}])_0\) along the principal
opens where \(b_j/b_i\) is invertible. The two overlap rings are both
\((B[b_i^{-1},b_j^{-1}])_0\), so the maps satisfy the cocycle identity.
This constructs \(\operatorname{Proj}B\). Gluing \((M[b_i^{-1}])_0\)
constructs \(\widetilde M\), and localization followed by a graded
direct summand proves exactness. For \(B=A[T_0,\ldots,T_n]\), the
\(i\)-th chart is \(\operatorname{Spec}A[T_j/T_i:j\ne i]\); on an
overlap its coordinates change by \((T_l/T_i)/(T_j/T_i)=T_l/T_j\). The
frame \(T_i^d\) glues the invertible sheaf \(\mathcal O(d)\), with
transition multiplier \((T_j/T_i)^d\). These constructions commute with
arbitrary change of \(A\), because localization, grading and their
transition maps do. A graded quotient
\(A[T_0,\ldots,T_n]\twoheadrightarrow B\) is surjective on every chart
ring and hence gives a closed immersion of Proj with the same twists.

The quotient universal property follows from the same charts. A
surjection \(\mathcal O_T^{n+1}\to L\) onto an invertible sheaf has
opens where its \(i\)-th image generates \(L\). On each such open the
ratios of the other images to that image define the coordinate map to
the \(i\)-th projective chart. The ratios agree on overlaps, so they
glue uniquely; the pullback of \(\mathcal O(1)\) is \(L\). Conversely
the coordinate sections give this quotient. Changing a free frame
therefore acts canonically on projective space and its tautological
quotient. For a finite locally free sheaf \(E\) on a general base, glue
these free constructions under its frame changes. The cocycle identity
and uniqueness of the quotient maps give
\(\mathbf P(E)=\operatorname{Proj}(\operatorname{Sym}E)\), representing
invertible quotients of pullbacks of \(E\), with
\(\pi^*E\twoheadrightarrow\mathcal O(1)\). For a finite-type \(E\) over
an affine base, a finite free surjection gives the same construction as
the closed Proj of its symmetric-algebra quotient. This also proves the
projective embedding used in the next lesson.

\subsection{1. The monomial section
complex}\label{the-monomial-section-complex}

Fix a ring \(A\), an integer \(n\ge0\), and the standard graded ring \[
S=A[T_0,\ldots,T_n],\qquad\deg T_i=1.
\] Write \(\mathbf P_A^n=\operatorname{Proj}S\). The twist
\(\mathcal O(d)\) is the sheaf of the graded module \(S(d)\), with
\(S(d)_m=S_{d+m}\). On a nonempty intersection of standard charts, \[
\Gamma(D_+(T_{i_0})\cap\cdots\cap D_+(T_{i_p}),\mathcal O(d))
=S[T_{i_0}^{-1},\ldots,T_{i_p}^{-1}]_d.
\] The subscript means homogeneous degree \(d\). Inverting the product
of these variables is equivalent to inverting each of them. These
intersections are affine; for example, the chart indexed by \(i\) has
coordinates \(T_j/T_i\). Thus affine-cover comparison identifies
cohomology with the ordered Čech complex \[
C^p(d)=\bigoplus_{i_0<\cdots<i_p}
S[T_{i_0}^{-1},\ldots,T_{i_p}^{-1}]_d.
\] It stops in degree \(n\). The differential is the alternating sum of
restrictions, with increasing indices and the sign of the deleted
index's position.

Each term is a free \(A\)-module on Laurent monomials. This remains true
when \(A\) has torsion or zero divisors: the monomials are formal basis
vectors of the polynomial localization. A weight is an integer vector
\(e=(e_0,\ldots,e_n)\), with total \(\sum e_i=d\). The monomial \(T^e\)
occurs in the factor indexed by a subset \(I\) exactly when \[
N(e)=\{i:e_i<0\}\subset I.
\] The differential preserves weights. The whole complex is therefore a
direct sum of complexes \(C^\bullet(e)\), each having a copy of \(A\)
for the nonempty subsets \(I\) containing \(N(e)\). Its differential
uses only zero maps and signed identity maps between these copies.

These are direct sums rather than products over the weights: every
Laurent polynomial has finite support, and only finitely many chart
subsets occur in a fixed degree. Direct sums are exact, so cohomology
can be computed separately at each weight.

\subsection{2. Negative exponents determine
cohomology}\label{negative-exponents-determine-cohomology}

\textbf{Lemma 2.1 (the three weight cases).} For a weight \(e\), the
cohomology of \(C^\bullet(e)\) is:

{
\begin{longtable}[]{@{}ll@{}}
\toprule\noalign{}
Negative indices & Surviving cohomology \\
\midrule\noalign{}
\endhead
\bottomrule\noalign{}
\endlastfoot
None & One copy of \(A T^e\) in degree zero \\
All \(n+1\) indices & One copy of \(A T^e\) in degree \(n\) \\
A nonempty proper subset & Zero in every degree \\
\end{longtable}
}

\textbf{Proof.} If every exponent is negative, only the subset
\(I=\{0,\ldots,n\}\) is allowed. Its copy of \(A\) is the whole complex,
in degree \(n\).

If no exponent is negative, every nonempty subset is allowed. This is
the ordinary simplex cochain complex, with \(A\) as its degree-zero
cohomology. More explicitly, augment it by a copy of \(A\) in degree
\(-1\), whose map is the diagonal. Inserting a fixed index gives a
contracting homotopy of this augmented complex. Thus the original
complex has the stated cohomology.

Now suppose \(N(e)\) is nonempty and proper. Choose \(j\notin N(e)\).
Regard an ordered cochain as an alternating function of distinct
indices, with value zero on any tuple whose set does not contain
\(N(e)\), and zero on repeated indices. Define \[
(hc)_{i_0\ldots i_{p-1}}=c_{j i_0\ldots i_{p-1}}.
\] Insertion is followed by reordering with its sign. This map is
well-defined: adjoining \(j\) does not change whether a subset contains
\(N(e)\). The identity \(dh+hd=1\) follows by pairing each deletion
other than the inserted \(j\) with its opposite-sign counterpart;
deletion of \(j\) contributes the original cochain. In degree zero the
would-be augmented term is absent because \(N(e)\ne\varnothing\), and
the same calculation still works: the value at the singleton \(\{j\}\)
is zero. This is a contraction of the unaugmented complex, so all its
cohomology is zero. \(\square\)

For a concrete example with three variables, take \(N(e)=\{0\}\). The
weight complex is \[
A\longrightarrow A^2\longrightarrow A,
\qquad a\longmapsto(-a,-a),\qquad(b,c)\longmapsto b-c.
\] The two middle entries correspond to \(\{0,1\}\) and \(\{0,2\}\). The
last map is surjective and its kernel is the diagonal, exactly the
preceding image. With \(N(e)=\{0,1\}\), the only terms are
\(A\xrightarrow{1}A\) in degrees one and two. With all three exponents
negative, only the last \(A\) survives. These small complexes illustrate
the general cancellation without dividing by an integer.

Put \(r=n+1\). For \(n\ge1\), let \[
L_d=\bigoplus_{\substack{e_i<0\text{ for all }i\\\sum e_i=d}} A T^e.
\] It is zero unless \(d\le-r\). It can also be written as the
homogeneous degree-\(d\) part of \[
\frac1{T_0\cdots T_n}A[T_0^{-1},\ldots,T_n^{-1}].
\] The expression is a module of reciprocal monomials, not a ring.

\textbf{Theorem 2.2 (cohomology of every twist).} If \(n\ge1\), then \[
H^q(\mathbf P_A^n,\mathcal O(d))=
\begin{cases}
S_d,&q=0,\ d\ge0,\\
L_d,&q=n,\ d\le-n-1,\\
0,&\text{otherwise}.
\end{cases}
\] All the nonzero modules are finite free over \(A\). For \(n=0\),
instead, \[
H^0(\mathbf P_A^0,\mathcal O(d))=A\quad\text{for every integer }d,
\qquad H^q=0\quad(q>0).
\]

\textbf{Proof.} For positive dimension, Lemma 2.1 and the weight direct
sum identify degree-zero classes with monomials having nonnegative
exponents, and top-degree classes with monomials having strictly
negative exponents. All other weights are contractible. The sum
condition gives the stated ranges. For fixed \(d\), there are only
finitely many vectors of nonnegative exponents summing to \(d\), or of
strictly negative exponents summing to \(d\), proving finite freeness.

In dimension zero there is one chart, and its degree-\(d\) section
module is \(A T_0^d\) for every integer \(d\). This is \(A\);
geometrically \(\mathbf P_A^0=\operatorname{Spec}A\), and the
trivializing generator of the twist is \(T_0^d\), including negative
powers. There are no higher Čech terms. \(\square\)

The separate zero-dimensional clause prevents a false negative-twist
vanishing statement when degree zero is also top degree. For \(n\ge1\),
all groups vanish in the entire interval \(-n\le d\le-1\). The
calculation is {[}Stacks, Tag 01XT{]}, with the two cohomological
degrees kept distinct when necessary.

\subsection{3. A pairing with no
factorials}\label{a-pairing-with-no-factorials}

The top group of \(\mathcal O(-r)\) is free of rank one. Fix its trace
by \[
\operatorname{tr}\left[\frac1{T_0\cdots T_n}\right]=1.
\] For \(m\ge0\), multiplication by global sections gives a bilinear
pairing \[
S_m\times H^n(\mathbf P_A^n,\mathcal O(-m-r))
\longrightarrow H^n(\mathbf P_A^n,\mathcal O(-r))
\xrightarrow{\operatorname{tr}}A.
\]

\textbf{Theorem 3.1 (perfect monomial pairing).} This pairing is
perfect: each module is the \(A\)-linear dual of the other.

\textbf{Proof.} Index the monomial basis of \(S_m\) by vectors
\(\alpha_i\ge0\) with \(\sum\alpha_i=m\). A corresponding basis of the
other group is \[
u_\alpha=T_0^{-\alpha_0-1}\cdots T_n^{-\alpha_n-1}.
\] The product \(T^\beta u_\alpha\) has exponent \(\beta_i-\alpha_i-1\).
If \(\beta=\alpha\), it is the trace generator. If the vectors differ
but have the same total, at least one \(\beta_i>\alpha_i\), so that
product has a nonnegative exponent and represents zero in top cohomology
by Lemma 2.1. The pairing matrix is therefore the identity. This proves
perfection over every ring, including rings with positive characteristic
or nilpotents. \(\square\)

The proof also applies to \(n=0\) using the generators of its one-chart
section modules. There is no appeal to integration and no normalization
by a factorial. This will be the local calculation behind Serre duality
later in the course.

\textbf{Proposition 3.2 (base change and polynomial multiplication).}
The preceding identifications commute with every ring map \(A\to B\),
including nonflat ones. If \(f\in S_s\), multiplication by \(f\) on top
cohomology is the dual of \[
S_{-d-s-r}\xrightarrow{\ f\ }S_{-d-r},
\] where negative homogeneous pieces are zero and \(n\ge1\).

\textbf{Proof.} The section complex over \(B\) is obtained by replacing
the coefficient ring in every monomial factor by \(B\). The weight
contractions use only signed identities, so they remain contractions
after any coefficient change. The surviving bases also change by
extension of scalars. The resulting isomorphism \[
H^q(\mathbf P_A^n,\mathcal O(d))\otimes_A B
\simeq H^q(\mathbf P_B^n,\mathcal O(d))
\] is the natural one induced by the chart section complexes. No general
nonflat base-change theorem is being assumed here.

For the multiplication statement, let \(u\) be a top class and \(g\) a
homogeneous polynomial of degree \(-d-s-r\). Associativity of
multiplication gives \[
\operatorname{tr}(g(fu))=\operatorname{tr}((fg)u).
\] The perfect pairing identifies this equality with the claimed dual
map. Products that leave the negative-exponent range are zero classes,
as already proved. \(\square\)

This functorial form agrees with {[}Stacks, Tag 01XV{]}. In particular,
the finite free modules in the calculation retain their ranks under
every specialization of the base ring.

\subsection{4. Higher direct images over any
base}\label{higher-direct-images-over-any-base}

Let \(\pi:\mathbf P_S^n\to S\), for an arbitrary scheme \(S\). In
positive dimension the formula becomes \[
R^q\pi_*\mathcal O(d)=
\begin{cases}
\operatorname{Sym}^d(\mathcal O_S^{r}),&q=0,\ d\ge0,\\
\mathcal H om(\operatorname{Sym}^{-d-r}(\mathcal O_S^{r}),\mathcal O_S),
&q=n,\ d\le-r,\\
0,&\text{otherwise}.
\end{cases}
\] The fixed ordered standard coordinates trivialize the determinant
line, which is why it does not appear in this particular formula. For
\(n=0\), the morphism is the identity and
\(\pi_*\mathcal O(d)=\mathcal O_S\) for every \(d\).

\textbf{Proof of the relative formula.} On an affine open
\(\operatorname{Spec}A\subset S\), the preceding lesson identifies
higher direct images with the sheaves of the modules
\(H^q(\mathbf P_A^n,\mathcal O(d))\). Theorem 2.2 computes these modules
and Theorem 3.1 gives their dual descriptions. Proposition 3.2 makes the
identifications compatible on distinguished opens; hence they glue
across arbitrary affine charts of \(S\). The same proposition shows that
the natural higher-direct-image base-change maps are isomorphisms for
every \(S'\to S\), by checking locally on source and target affines.
This supplies the relative proof of {[}Stacks, Tag 01XW{]}. \(\square\)

\subsection{5. The determinant line of a projective
bundle}\label{the-determinant-line-of-a-projective-bundle}

Let \(\mathcal E\) be finite locally free of constant positive rank
\(r\) on \(S\). We use the quotient convention \[
\mathbf P(\mathcal E)=\operatorname{Proj}_S\operatorname{Sym}(\mathcal E),
\qquad\pi^*\mathcal E\twoheadrightarrow\mathcal O(1).
\] Thus \(\pi_*\mathcal O(1)=\mathcal E\) when \(r\ge2\); the
corresponding assertion for rank one follows below. This convention is
important: defining projective space from \(\mathcal E^\vee\) would
change the displayed bundles. The construction and tautological quotient
were proved in the chart construction above. Write
\(D=\det\mathcal E=\bigwedge^r\mathcal E\).

\textbf{Lemma 5.1 (canonical top trace).} For \(r\ge2\), there is a
canonical isomorphism \[
D\otimes R^{r-1}\pi_*\mathcal O(-r)\simeq\mathcal O_S.
\] In a local ordered frame it takes the frame's exterior product
tensored with the reciprocal-coordinate class to one.

\textbf{Proof.} Twist the tautological quotient down to obtain a
surjection \[
F=\pi^*\mathcal E\otimes\mathcal O(-1)\longrightarrow\mathcal O.
\] Its Koszul differential on \(\bigwedge^j F\) is contraction with this
functional. The resulting sequence \[
0\to\pi^*D\otimes\mathcal O(-r)\to\cdots\to
\pi^*\mathcal E\otimes\mathcal O(-1)\to\mathcal O\to0
\] is exact. Indeed locally a vector \(v\in F\) maps to one; exterior
multiplication by \(v\) is a contracting homotopy because \[
\iota(v\wedge w)+v\wedge\iota(w)=w.
\] This also shows that its images are locally direct summands, so the
sequence remains exact under any base change.

For \(1\le j\le r-1\), every higher direct image, including degree zero,
of \(\pi^*(\bigwedge^j\mathcal E)\otimes\mathcal O(-j)\) is zero. This
can be checked on a trivializing open: it is a finite direct sum of
copies of the acyclic twist \(\mathcal O(-j)\), by Section 4. Breaking
the Koszul sequence into short exact sequences of its images and
composing their connecting maps therefore gives a canonical isomorphism
\[
\mathcal O_S=\pi_*\mathcal O
\xrightarrow{\delta}R^{r-1}\pi_*(\pi^*D\otimes\mathcal O(-r)).
\] For a finite locally free base module \(G\), the natural projection
map identifies \(R^q\pi_*(\pi^*G\otimes\mathcal F)\) with
\(G\otimes R^q\pi_*\mathcal F\): on a trivializing open it is just
commutation of cohomology with a finite direct sum. Applying this to
\(D\) and inverting \(\delta\) gives the trace in the statement.

To check its normalization, choose a frame \(e_0,\ldots,e_{r-1}\), with
corresponding coordinates \(T_i\). On the \(i\)-th chart put
\(v_i=e_i/T_i\) in \(F\); its Koszul contraction is one. The Čech
cochains \[
v_{i_0}\wedge\cdots\wedge v_{i_p}
\] are successive lifts in the connecting-map computation: their Koszul
differential is the alternating deletion sum of the preceding cochain.
Thus \(\delta(1)\) is represented on the full intersection by \[
\frac{e_0\wedge\cdots\wedge e_{r-1}}{T_0\cdots T_{r-1}}.
\] This proves the stated normalization, and the construction from the
tautological quotient proves independence of the frame. \(\square\)

This is the canonical Koszul argument of {[}Stacks, Tag 01XX{]}. It
determines the determinant factor without assuming that every invertible
matrix over the base ring has a chosen elementary factorization.

\textbf{Theorem 5.2 (cohomology of a projective bundle).} If \(r\ge2\),
then \[
R^q\pi_*\mathcal O(d)=
\begin{cases}
\operatorname{Sym}^d\mathcal E,&q=0,\ d\ge0,\\
(\operatorname{Sym}^{-d-r}\mathcal E\otimes D)^\vee,
&q=r-1,\ d\le-r,\\
0,&\text{otherwise}.
\end{cases}
\] These are canonical identifications compatible with isomorphisms of
\(\mathcal E\) and arbitrary base change. For rank one,
\(\mathbf P(\mathcal E)=S\), \(\mathcal O(d)=\mathcal E^{\otimes d}\),
and all higher direct images vanish; negative powers mean powers of the
dual line bundle.

\textbf{Proof.} Trivializing \(\mathcal E\) reduces vanishing to Section
4. The natural map
\(\operatorname{Sym}^d\mathcal E\to\pi_*\mathcal O(d)\) is an
isomorphism for \(d\ge0\) by the same local calculation. For \(d=-m-r\),
multiplication and Lemma 5.1 give the canonical pairing \[
\operatorname{Sym}^m\mathcal E\otimes
R^{r-1}\pi_*\mathcal O(-m-r)\longrightarrow D^\vee.
\] In any ordered local frame, Lemma 5.1 identifies it with the perfect
pairing of Theorem 3.1. It consequently induces an isomorphism to
\((\operatorname{Sym}^m\mathcal E\otimes D)^\vee\). Perfection is local,
so this identifies the global bundles. All maps used are natural in the
tautological quotient; local monomial base change and the normalized
trace show compatibility with arbitrary base change. This proves the
positive-rank formula.

In rank one the tautological quotient of the pulled-back line is an
isomorphism of line bundles. The quotient-line functor is represented by
\(S\) itself: a surjection from a line to a line is an isomorphism,
leaving one isomorphism class of quotient at every base change. Hence
\(\mathbf P(\mathcal E)=S\) and \(\mathcal O(1)=\mathcal E\), giving the
separate clause. \(\square\)

The dual on a symmetric power should be left as the dual of that
symmetric power. Over general rings one must not silently replace it by
the symmetric power of the dual using a characteristic-zero
symmetrization argument. The coordinate pairing above works without such
a replacement.

For a finite locally free sheaf of varying rank, apply the theorem on
its open rank loci. A rank-zero locus has empty projective bundle, since
its symmetric algebra has no positive-degree elements; all direct images
there are zero. Thus the preceding statements cover arbitrary finite
locally free \(\mathcal E\), including rank one and rank zero.

\subsection{6. Ranks and Euler
characteristics}\label{ranks-and-euler-characteristics}

On \(\mathbf P_A^1\), the top group of \(\mathcal O(-2)\) is \(A\),
generated by \(1/(T_0T_1)\). With \(t=T_1/T_0\) and the frame
\(T_0^{-2}\) on the first chart, its coefficient is \(t^{-1}\), agreeing
with the two-chart calculation in \emph{Čech cohomology}.

More generally, on \(\mathbf P_k^n\) over a field, define \[
\chi(\mathcal O(d))=\sum_q(-1)^q\dim_k H^q(\mathbf P_k^n,\mathcal O(d)).
\] For \(d\ge0\), stars and bars counts \(\binom{d+n}{n}\) nonnegative
exponent vectors. For \(d\le-n-1\), write \(e_i=-\alpha_i-1\); then
\(\sum\alpha_i=-d-n-1\), giving \(\binom{-d-1}{n}\) top basis vectors.
Their Euler contribution has sign \((-1)^n\). In the intervening range
every group is zero. All three cases are expressed by the same
polynomial: \[
\chi(\mathcal O(d))=\binom{d+n}{n}
=\frac{(d+1)\cdots(d+n)}{n!}\qquad(d\in\mathbf Z).
\] Here the binomial is the generalized binomial polynomial, not a
convention that sets negative upper arguments to zero. In the negative
range reversing the signs of its \(n\) factors gives exactly
\((-1)^n\binom{-d-1}{n}\); in the intervening range a factor is zero.
For \(n=0\), the empty product is one, agreeing with \(\chi=1\) for
every twist.

The same numbers count the ranks of the finite free cohomology modules
over an arbitrary ring. This is a statement about integral ranks, not
division by \(n!\) inside the base ring.

\subsection{7. Exercises with solutions}\label{exercises-with-solutions}

\textbf{Exercise 7.1 (easy: every twist on a line).} Compute
\(H^q(\mathbf P_A^1,\mathcal O(d))\) for every integer \(d\), giving
bases whenever it is nonzero.

\textbf{Solution.} For \(d\ge0\), the degree-zero basis is
\(T_0^{d-j}T_1^j\), with \(0\le j\le d\), and the first group is zero.
For \(d=-1\), both groups vanish. For \(d\le-2\), the degree-zero group
vanishes and the first group has basis \[
T_0^{-i}T_1^{-j},\qquad i,j\ge1,\quad i+j=-d.
\] There are \(-d-1\) such vectors. Degrees greater than one vanish.
These conclusions follow from the two-chart complex: its first
cohomology kills precisely Laurent monomials with at least one
nonnegative homogeneous exponent. The groups are free \(A\)-modules on
the listed bases, regardless of torsion in \(A\).

\textbf{Exercise 7.2 (easy: the dual bases).} For \(m=3\) on
\(\mathbf P_A^1\), write the two dual bases in Theorem 3.1 and check
every pairing entry.

\textbf{Solution.} The polynomial basis is
\(T_0^3,T_0^2T_1,T_0T_1^2,T_1^3\). The reciprocal basis of
\(H^1(\mathcal O(-5))\) is \[
T_0^{-4}T_1^{-1},\quad T_0^{-3}T_1^{-2},\quad
T_0^{-2}T_1^{-3},\quad T_0^{-1}T_1^{-4}.
\] Matching entries multiply to \(T_0^{-1}T_1^{-1}\), with trace one.
For unequal entries the two exponent sums still total \(-2\); one
exponent is nonnegative, so the product is a zero class. Thus every
off-diagonal entry is zero and every diagonal entry is one. This
identity matrix is invertible over any \(A\), not merely over fields.

\textbf{Exercise 7.3 (medium: the projective plane).} Compute all groups
on \(\mathbf P_k^2\) and determine the rank of \(H^2(\mathcal O(-5))\).
Explain why weights with only one or two negative exponents contribute
nothing.

\textbf{Solution.} For \(d\ge0\), only \(H^0\) is nonzero, with
dimension \(\binom{d+2}{2}\). For \(d=-1,-2\), every group vanishes. For
\(d\le-3\), only \(H^2\) is nonzero, with basis
\(T_0^{-i}T_1^{-j}T_2^{-\ell}\), where \(i,j,\ell\ge1\) and
\(i+j+\ell=-d\). Its dimension is \(\binom{-d-1}{2}\). There is no first
cohomology for any twist, and no groups in degrees greater than two.

For \(d=-5\), the positive triples are the three permutations of
\((3,1,1)\) and the three permutations of \((2,2,1)\); there are six,
agreeing with \(\binom42=6\). With one negative index the weight complex
is \(A\to A^2\to A\), with maps \((-1,-1)\) and \((1,-1)\), and is
exact. With two negative indices it is a signed identity \(A\to A\).
These are the mixed-sign cancellations that eliminate intermediate
classes.

\textbf{Exercise 7.4 (medium: the binomial polynomial).} Verify the
Euler-characteristic formula for \(n=2\), including \(d=-1,-2,-3,-4\).
What changes for projective dimension zero?

\textbf{Solution.} The polynomial is \((d+1)(d+2)/2\). At \(-1,-2\), it
is zero, agreeing with total vanishing. At \(-3,-4\), it is one and
three, the dimensions of the top groups; their Euler sign is positive
because the top degree is two. For \(d\ge0\), it counts polynomial
monomials. For arbitrary negative \(d\le-3\), it equals
\(\binom{-d-1}{2}\), so these cases cover every integer. In dimension
zero every twist is a trivial line on \(\operatorname{Spec}k\), and
\(\chi=1\) for all \(d\). The generalized binomial \(\binom d0\) and its
empty-product interpretation are also one. Treating a negative upper
binomial argument as automatic zero would give the wrong answer.

\textbf{Exercise 7.5 (challenging: a product of projective lines).} Over
any ring \(A\), compute the cohomology of \(\mathcal O(a,b)\) on
\(\mathbf P_A^1\times_A\mathbf P_A^1\). Give a proof that justifies a
tensor-product computation without an unproved Künneth assumption.

\textbf{Solution.} Put \(V_0(d)=H^0(\mathbf P_A^1,\mathcal O(d))\) and
\(V_1(d)=H^1(\mathbf P_A^1,\mathcal O(d))\), with the free bases of
Exercise 7.1. The answer is \[
\begin{aligned}
H^0&=V_0(a)\otimes_A V_0(b),\\
H^1&=(V_0(a)\otimes_A V_1(b))\oplus(V_1(a)\otimes_A V_0(b)),\\
H^2&=V_1(a)\otimes_A V_1(b),
\end{aligned}
\] with all later groups zero. Tensoring the two monomial bases gives
explicit bases. The ranks are obtained using
\(\operatorname{rank}V_0(d)=\max(d+1,0)\) and
\(\operatorname{rank}V_1(d)=\max(-d-1,0)\).

Here is the justification. Use the standard two-chart cover in each
factor. The two-direction Čech complex on the product is the tensor
product over \(A\) of their section complexes; each bi-intersection is
affine and its section monomials are products of the two factor
monomials. Its total complex computes sheaf cohomology by the same
acyclic-cover double-complex proof used earlier: on stalks the augmented
complexes in both cover directions have an insertion contraction, and on
the bi-intersections positive quasi-coherent cohomology vanishes. The
total differential is \(d\otimes1+(-1)^p1\otimes d\) in bidegree
\((p,q)\).

Each two-term factor complex splits into its cohomology and contractible
pieces already over \(A\). To see this explicitly, write it in the first
chart frame as \[
A[t]\oplus t^d A[t^{-1}]\longrightarrow A[t,t^{-1}],
\qquad(u,v)\longmapsto v-u.
\] At an exponent present in both source modules the map is
\(A^2\to A\), \((u,v)\mapsto v-u\); the diagonal is the degree-zero
class and \((0,1)\) splits the map. At an exponent in exactly one source
module the map is a signed identity and is contractible. At an exponent
in neither source module there is just a copy of \(A\) in degree one.
These decompositions exhaust every exponent. Tensoring a contraction
with the other complex, using the total-complex signs, still gives a
contraction. Hence only the tensor products of the two cohomology
summands survive. This proves the displayed formulas without any
flatness or field assumption.

\textbf{Exercise 7.6 (medium: detect the determinant convention).} Let
\(\mathcal E=\mathcal L\oplus\mathcal M\), for line bundles on \(S\).
Compute \(R^1\pi_*\mathcal O(-3)\). Then compare with the rank-one
bundle \(\mathbf P(\mathcal L)\).

\textbf{Solution.} Here \(r=2\), \(-d-r=1\), and
\(\det\mathcal E=\mathcal L\otimes\mathcal M\). Theorem 5.2 gives \[
R^1\pi_*\mathcal O(-3)
=(\mathcal L^{\otimes2}\otimes\mathcal M)^\vee
\oplus(\mathcal L\otimes\mathcal M^{\otimes2})^\vee.
\] The zeroth direct image is zero and the groups in degrees at least
two vanish. If both lines are trivial, the displayed bundle is
\(\mathcal O_S^{\oplus2}\), matching the two reciprocal monomials of
degree \(-3\) on a projective line. In rank one, instead, the projective
bundle is \(S\) itself, \(\pi_*\mathcal O(-3)=\mathcal L^{-3}\), and the
first direct image is zero. Applying the positive-dimensional
negative-twist rule without the rank-one clause would lose this nonzero
zeroth direct image.

\subsection{References and construction
prerequisites}\label{references-and-construction-prerequisites}

\begin{itemize}
\tightlist
\item
  \textbf{{[}Stacks{]}} The Stacks project authors, \emph{The Stacks
  project}, in its AI Integrated Stacks Project edition. Cohomology over
  a ring:
  \href{https://kokunoyumeto.github.io/stacks-zh-hans-cn/en/coherent.html\#coherent-lemma-cohomology-projective-space-over-ring}{Tag
  01XT}; functoriality and the monomial pairing:
  \href{https://kokunoyumeto.github.io/stacks-zh-hans-cn/en/coherent.html\#coherent-lemma-identify-functorially}{Tag
  01XV}; arbitrary base:
  \href{https://kokunoyumeto.github.io/stacks-zh-hans-cn/en/coherent.html\#coherent-lemma-cohomology-projective-space-over-base}{Tag
  01XW}; the canonical projective-bundle argument:
  \href{https://kokunoyumeto.github.io/stacks-zh-hans-cn/en/coherent.html\#coherent-lemma-cohomology-projective-bundle}{Tag
  01XX}. Complete arguments for these computations and the
  dimension-zero cases are given above.
\item
  Construction prerequisites: Proj and its twists
  \href{https://kokunoyumeto.github.io/stacks-zh-hans-cn/en/constructions.html\#constructions-section-proj}{Tag
  01M3}, and the quotient convention and tautological line for
  projective bundles
  \href{https://kokunoyumeto.github.io/stacks-zh-hans-cn/en/constructions.html\#constructions-section-projective-bundle}{Tag
  01OA}. The chart construction at the start of this lesson proves the
  degree-one Proj, twists, quotient convention and base-change
  statements actually used here and in the next lesson. These open
  reference treatments retain the GNU Free Documentation License 1.2.
  The prose, proofs, examples and solutions written here are independent
  CC0 content.
\item
  AI Integrated Stacks Project contains AI-proposed corrections and
  additions, and is not reviewed by maintainers of the
  \href{https://stacks.math.columbia.edu/}{official Stacks project}. All
  calculations use the stated quotient convention and preserve the
  actual ring generality.
\end{itemize}
\end{document}
